% LAB BOREDOM — module L0, Introduction. GUIDED DRAFTING (desktop recipe). Session 2026-08-21.
% Spec: outline v8 §L0 AS AMENDED by author rulings of 2026-08-21:
%   ★ THESIS RESTRUCTURED (author): film-specific framing COMPLETELY REPLACED by the overall
%     point — boredom/engagement identifiable in physiological data and subjective surveys.
%     The nine divergent episodes are ONE result beside the closure/gaze-deviation
%     significance, NOT the thesis's primary point. The MAIN point: Heidegger's phenomenology
%     applied to the data generates insight neither pure phenomenology nor pure statistical
%     analysis can provide.
%   ★ Resolution move recast (approved with plan): episode as unit + criterion analysis as
%     the experiment — the old two-questions (separation/agreement) frame is void.
%   Structure: hook > context > disruption > resolution > signpost > thesis LAST (the
%     author's locked intro structure). Target 9 ¶¶: B11 ¶¶1–3 · B12 ¶¶4–6 · B13 ¶¶7–9.
% Standing rulings carried (see L4–L7 headers): never "laboratory"; data-as-agent; no
%   deployment mention anywhere in L0; episode-only, film names at most as identifiers; no
%   "ones"; subjects by number, ungendered; tree-convention figures; percentages numerals.
%
% BATCH 11 = ¶¶1–3 (hook + context). Drafted 2026-08-21, awaiting author review.
% Verification ledger, batch 11:
%   ¶1 hook facts: S07 CLC boredom 6 / engagement 7; S07 INT boredom 7 / engagement 6
%     (survey-episodes.csv, verified this session); survey administered VERBALLY by the
%     proctor in the ~5-minute breaks; 17-minute films; separate boredom and engagement
%     items (L2-established). The 4-episodes/3-participants count HELD for later modules.
%   ¶2 context facts: Elpidorou as the self-citing author bridging the two programs + the
%     single critical footnote (Hughes), twenty-four years — L1 ¶1's established claims,
%     compressed, not re-argued.
%   ¶3 lineage facts: clinical-immersion pedagogy, scarcity-of-clinical-access warrant
%     (grp-program-lineage, L1 ¶8a); definitional adoption Eastwood -> Fahlman -> P1
%     (L1 ¶8b); definition PARAPHRASED per D-23 (verbatim renders once, in L1 ¶10).
% BATCH 11 REV 1 (2026-08-21, author notes): ¶1 structure NAMED (an episode that genuinely
%   engages and genuinely leaves the viewer empty at once). ¶2 final sentence DELETED
%   (24-years/self-citing detail belongs to L1's review; the intro's review stays less
%   specific — the section carries two separate literature reviews by design). ¶3 approved
%   CONDITIONAL on ¶4 cashing out the question, the reasonableness, and the consequence —
%   B12 opens with exactly that.
%
% BATCH 12 = ¶¶4–6 (disruption + resolution). Drafted 2026-08-21, awaiting author review.
% Verification ledger, batch 12:
%   ¶4 the D-19 preview at intro altitude (the full argument lives in L1; not made twice at
%     full strength): reasonable = validated state-boredom construct with instruments behind
%     it, stimulus selection following it; consequential = two of the three forms deny
%     environmental attribution (second: no culprit findable; third: impersonal), so the
%     survey is a first-form instrument.
%   ¶5 disruption facts (established): EEG headline from one clean participant, full
%     negative at 8 (0 of 5 measures); published gaze feature rate-bound (defined at the
%     sensor's 120 Hz; Round 1 telemetry ~3 Hz); the design-implied criterion analysis never
%     computed before this reanalysis.
%   ¶6 resolution facts: both rounds reprocessed from raw; 12 participants x 3 = 36
%     episodes; episode as unit; within-participant scoring; the two admitted eye measures;
%     whole-survey scoring; genre statement per spec 0a (two moves, seam visible).
% BATCH 12 REV 1 (2026-08-21, author notes): ¶4 APPROVED. ¶5 the never-computed claim
%   CORRECTED per author — the criterion analysis existed in part (subsets computed,
%   qualitative assertions for the rest, per D-16); what had never been run is the full
%   twelve-participant, thirty-six-episode version. ¶6 "and the seam ... they meet." DELETED.
%
% BATCH 13 = ¶¶7–8 (signpost + thesis, thesis LAST). Drafted 2026-08-21, awaiting review.
%   L0 lands at 8 ¶¶ total (one under the ~9 target; flagged to author).
%   ¶7 the seven-part roadmap, functional enumeration on the M0.4 model; the seventh part
%     described as closing Part III WITHOUT naming the classroom study (deployment barred
%     from L0; the roadmap only points at the close).
%   ¶8 the thesis per the author's amended structure: (1) boredom/engagement identifiable
%     in physiological data and subjective surveys, each identifying the other; (2) the
%     divergence as ONE result among the correlations, readable because agreement is the
%     rule; (3) MAIN: the Heideggerian interpretive frame generates insight neither pure
%     phenomenology nor pure statistical analysis provides. "demonstrate" to the last line.
% BATCH 13 REV 1 (2026-08-21, author notes):
%   Author clarification resolved: L0 introduces SECTION 2 ONLY (Part III has its own intro,
%     III-0; Section 1 has M0). ¶¶1–6 REVIEWED on author order for Section 1 over-discussion:
%     CLEAN — no classroom-study material anywhere in them (¶3's "third home" is the
%     clinical-immersion VR lineage of the films, not the VLE study).
%   P7-1: roadmap's final sentence made specific — states what the close will demonstrate
%     (survey-only attests first form; the experiment's two kinds of evidence reach the
%     second; the Part hands forward a frame that survived contact with evidence in one
%     study and did work on evidence in the other). One-sentence naming of the classroom
%     study in the SIGNPOST is sanctioned by the author's own request for specificity; the
%     deployment bar otherwise holds throughout L0.
%   P8 (thesis) APPROVED as drafted.
% BATCH 13 REV 2 (2026-08-21, author note): the hands-forward clause anchored in specifics —
%   the frame described a cohort's collected responses without strain (classroom study) and
%   read out of the recorded eyes and the ratings a form of boredom the survey's own
%   definition could not name (this experiment).
% ★ ¶7 APPROVED 2026-08-21 (rev 2). ★★ MODULE L0 COMPLETE — 8 ¶¶.
% ★★★ SECTION BODY COMPLETE 2026-08-21: L0–L7 all drafted and author-approved. The appendix
%     deferral ("until the body is finished") now LIFTS. NOT committed — sign-off required.

In the physiological experiment this section analyzes, a participant recorded as S07 sat
through a seventeen-minute film in a headset, took a short break, and answered a survey the
proctor read aloud. Asked how bored, S07 answered 6 of 9. Asked how engaged, S07 answered 7.
On another of S07's three episodes that afternoon, the two answers traded places: bored 7,
engaged 6. A survey built on the assumption that boredom and engagement oppose each other
would have folded each pair into one middling number. This survey asked the two questions
separately, and the instrument that allowed two answers got two. Episodes like
S07's---affirmed as boring and engaging in the same breath---recur in this data, and they are
neither error nor indecision. They are the first sighting of the structure this section
traces from its findings to its close: an episode that genuinely engages its viewer and
genuinely leaves that viewer empty, at the same time.

Boredom has three scholarly homes, and the two oldest have barely spoken. The first is
phenomenological: Heidegger's lecture course of 1929--30 made boredom a fundamental
attunement and asked what it discloses about the being who suffers it, and a tradition
running from pedagogy to film theory has worked in its wake. The second is empirical:
psychology and psychophysiology made boredom a state to be scaled, correlated, and lately
classified from signals, asking not what boredom discloses but what it looks like in a
measurement. The two traditions share the word and almost nothing else.

The studies that produced this data occupy a third home. They belong to clinical-immersion
pedagogy---virtual reality built to give biomedical-engineering students the hospital
experience that hospitals cannot scale---and they took their survey's definition of boredom
from the second home: an aversive state of wanting, but being unable, to engage in
satisfying activity, a state attributed to one's environment. The choice was reasonable, and
it was consequential. What a survey built under that definition can find, and what it must
miss, is a question the studies themselves never had reason to ask. It is where this section
begins.

The question has a precise form, and the two halves of this section's claim both hang on it.
The definition was a reasonable adoption because it is what a measuring study needs: a
construct built for measurement, with validated instruments behind it and a stimulus logic
that follows from it---material that is monotonous, redundant, or meaningless raises the
likelihood of the state. The adoption was consequential because of what the definition
requires: a bored person who locates the cause of their state in the environment. Heidegger's
analysis of boredom, the frame this dissertation carries, distinguishes three forms of the
attunement, and only the first has that shape. In the second form there is no culprit to
locate---the situation genuinely occupies, and the emptiness forms in the person---and in the
third the question of a culprit does not arise at all. A survey operationalizing the adopted
definition is therefore an instrument of the first form. Whatever else is present in an
episode, such a survey can register it only in first-form terms---or as answers it has no way
to interpret.

That question had never been put to this data, and parts of the data had gone unread. The
published headline from these recordings rested on the electroencephalogram of one clean
participant; reprocessed across all eight who wore the cap, the EEG is a full negative, with
none of its five derived measures distinguishing what participants reported. The published
gaze feature is defined at the sensor's full sampling rate, which the first round's telemetry,
logging at roughly 3 Hz, never delivered. Deepest of the three: the analysis the design itself
implied---every recorded measure scored against what the participant reported, episode by
episode---existed only in part, computed for subsets of participants and asserted
qualitatively for the rest; across all twelve participants and all thirty-six episodes it had
never been run.

This section computes that analysis, and everything else follows from how. Both rounds of
recordings were reprocessed from the raw files: twelve participants, three episodes each,
thirty-six episodes. The episode is the unit of analysis throughout---each episode carries
what the eye tracker recorded during it and everything its participant reported about
it---and every claim is within-participant, each person scored against their own
three-episode baseline, so that no one's habits of rating or of looking can stand in for
anyone else's. Two eye measures survive reprocessing with their validity intact, the
proportion of an episode the eyes spent closed and the distance the gaze wandered from that
participant's own habitual direction, and they are scored against the whole survey: boredom,
engagement, the fight against sleep, the minute boredom arrived, felt duration, and fatigue
before and after. Every part of this section then makes one of two moves---here is what the
measurements show, or here is what those measurements mean once Heidegger's account of
boredom is the operative one.

Seven parts follow. The first places this study among its three literatures and shows why
the phenomenological and the empirical traditions have to be brought together over this
data. The second lays out the apparatus and the data: the cohorts, the headset and its
sensors, the survey and its definition, and the rules that determine what may be compared
with what. The third presents the findings of the reanalysis. The fourth is the analysis
proper, in which the chain this dissertation has carried since its first part is walked
through the episodes and Heidegger's three forms of boredom are given their instrumental
shape. The fifth states what the data demonstrate in the form of an instrument: an index of
divergence and the three states it sorts. The sixth names the limits, each stated and
stopped. The seventh closes the Part: it joins this experiment to the VLE experiment that
precedes it, demonstrates that a survey alone can attest only the first form of boredom
while the experiment's two kinds of evidence reach the second, and states what Part III
hands forward: a Heideggerian frame that in the VLE experiment described a cohort's
collected responses without strain, and in this experiment read out of the recorded eyes and
the ratings a form of boredom the survey's own definition could not name.

The thesis has three connected parts, and the last is the main one. First, boredom and
engagement are identifiable in physiological data and in subjective surveys, and the two
kinds of evidence identify each other: across the thirty-six episodes, how far the gaze
wandered from a participant's own habitual direction and how much of an episode the eyes
spent closed both track what that participant reported, significantly and in the predicted
directions, while the survey's own items---the fight against sleep, fatigue, the minute
boredom arrived---converge on the same state. Second, the results include a structure as
well as correlations: in nine episodes the two eye measures ran engaged while the
participant's verdict said bored, a divergence that is readable as evidence precisely
because agreement is the rule around it. Third---and this is the main claim---Heidegger's
phenomenology of boredom, brought to this data as an interpretive frame, generates insight
that neither pure phenomenology nor pure statistical analysis provides: it reads the
divergence as the structural grammar of a form of boredom the collecting definition could
not name; it locates the survey and the eye tracker on either side of an involuntary act of
judgment, one recording before the verdict and one collecting it after; and it turns a set
of correlations into a description of what an episode of mandated watching is for the
person held inside it. Phenomenology without the data has no measurements; statistical
analysis without the frame has no way to say what its exceptions mean. Together they
demonstrate more than either could alone.
